\hypertarget{focal-plane-opa-beam-propagation-design-and-geometry-optimisation}{%
\chapter{Focal-plane OPA: beam-propagation design and geometry
optimisation}\label{fpr_design-focal-plane-opa-beam-propagation-design-and-geometry-optimisation}}

\begin{quote}
Follow-up: the power figures below omit the cross axis and the emitter
facet loss (about 21\% into air). With both included and a collimator
glued to the facets, a 1.2 mK tweezer needs about 57 mW instead of 43 mW
(\texttt{docs/stacked\_relay.md}, \texttt{docs/collimator.md}).
\end{quote}

Follow-up to \texttt{docs/independent\_steering.md} (route E,
sliding-window focal-plane OPA). Code: \texttt{picsim/bpm.py},
\texttt{picsim/opa/fpr\_layout.py}, \texttt{fpr\_bpm.py},
\texttt{star\_rs.py}, \texttt{focal\_plane.py}; scripts
\texttt{m10\_fpr.py}, \texttt{m10\_fpr\_meep.py},
\texttt{m11\_optimise.py}; results \texttt{results/m10\_*},
\texttt{results/m11\_*}. Simulated values are valid in the regime
simulated and should be reviewed before they are quoted. Design choices
below are made on physical and fabrication grounds, not to match the
original proposal\textquotesingle s numbers.

\hypertarget{1-result}{%
\section{Result}\label{fpr_design-1-result}}

A dense-port star coupler designed by wide-angle beam propagation, with
a Rowland-type port arc, 1.5x oversampled arms and wide emitter facets,
gives a focal-plane OPA in which 20 tweezers move and switch
independently along a row, at 1.5 times the power efficiency of the
first design. Verified by beam propagation of the full 128-port, 192-arm
structure (calibrated effective-index model, least-squares port
synthesis):

{\small
\begin{longtable}[]{@{}>{\raggedright\arraybackslash}p{0.485\dimexpr(\linewidth-6\tabcolsep)\relax}>{\raggedright\arraybackslash}p{0.330\dimexpr(\linewidth-6\tabcolsep)\relax}>{\raggedright\arraybackslash}p{0.184\dimexpr(\linewidth-6\tabcolsep)\relax}@{}}
\toprule\noalign{}
quantity & first design (128 x 128, confocal) & optimised design D1 \\
\midrule\noalign{}
\endhead
\bottomrule\noalign{}
\endlastfoot
field-averaged peak per unit port power (relative) & 1.00 & 1.50 \\
edge (60 um) / centre peak & 0.52 & 0.70 \\
laser power per 1.2 mK tweezer, field mean (fibre) & \textasciitilde65
mW & \textasciitilde43 mW \\
20 tweezers over +/-60 um & \textasciitilde1.3 W & \textasciitilde0.86
W \\
residual arm-phase aberration, worst port & up to 1.3 rad pk-pk &
\textless= 0.08 rad rms \\
others\textquotesingle{} change when one of 20 moves by up to 4 um &
\textless= 3.5 nm, \textless= 0.12\% & \textless= 3.5 nm, \textless=
0.11\% \\
residual at a switched-off site & 1.3e-6 & 1.6e-6 \\
static peak spread of 20 tweezers & 0.56\% & 0.52\% \\
spot FWHM along the row & 1.71 um & 1.71 um \\
\end{longtable}
}

Laser power uses the loss chain of \texttt{m9\_E\_power} (facet -0.37
dB, 7-stage tunable tree -0.39 dB, 1 cm at 0.3 dB/cm, meta-objective
transmission 0.7 assumed) and a Gaussian 1.04 um profile across the row.

\textbf{Design D1.} 128 input ports and 192 arms, both 1.7 um wide at
2.0 um pitch at the slab (0.3 um gaps); slab radius R = 1329 um; ports
on an arc of radius 0.555 R through the slab vertex, pointing at the
arm-array centre; arms on an arc of radius R; 30 um tapers to 0.6 um
guides; 192 end-fire emitters at 1.5 um pitch with 1.2 um facets (facet
taper \textless= 15 um); static phase trims on all arms and ports; port
amplitudes from least-squares synthesis on the calibrated response.

\hypertarget{2-method-and-its-validation}{%
\section{Method and its
validation}\label{fpr_design-2-method-and-its-validation}}

\textbf{Vertical dimension.} The 400 nm film is uniform and fully
etched, so every guide and the slab share one vertical mode, and the
lateral effective-index model carries the star-coupler physics (standard
AWG practice). The plain effective-index model overestimates lateral
confinement: its coupling lengths are 31 to 36\% short of the 3D
vector-FEM values. We calibrate the lateral core and cladding indices
(1.8423, 1.3551) to vector-FEM effective indices of 0.6, 1.0 and 1.7 um
guides (residuals \textless= 4.3e-4) and to the coupling length of the
dense ports (1.7 um, 0.3 um gap: 622 vs 625 um). Held-out checks: 1.4 um
guides at 0.49 um gap, 1373 vs 1330 um (+3\%); 0.6 um guides at 1.0 um
gap, 2287 vs 1777 um (+29\%; narrow, widely spaced guides couple
negligibly in the design, so this residual has little effect). A second
calibration that fits both coupling lengths degrades the dense-port
value by 10\% and is kept as a bound.

\textbf{Propagation.} 2D wide-angle (Pade 1,1) Crank-Nicolson BPM with
the TM-type lateral operator that matches the chip TE mode; all ports
propagated together. Unit tests: power conserved to 1e-3 over 300 um;
phase index error \textless{} 3e-4; half a coupling length transfers
half the power. Grid: dx = 0.04 um, dz = 0.125 um (capture within 0.5\%
and arm-field fidelity \textgreater= 0.992 against dz = 0.0625 um).

\textbf{Against 2D FDTD.} On a reduced 8-port, 8-arm coupler with the
same angular range, MEEP (resolution 20/um) and the BPM agree with
arm-field fidelity 0.997 (centre port), 0.985 (middle) and 0.93 (edge
port), and capture within 2 to 5\%. At FDTD resolution 30/um the
edge-port fidelity rises to 0.981 (rms phase difference 0.13 rad,
capture 0.369 vs 0.371), and the FDTD result itself changes by fidelity
0.98 between the two resolutions; hence the residual difference is
dominated by FDTD discretisation rather than by the BPM.

\textbf{Exact-path (Huygens) model.} For fast optimisation, a 2D
Rayleigh-Sommerfeld model with exact path lengths (\texttt{star\_rs.py})
reproduces the full 128 x 128 BPM to within 1\% in capture and with
fidelity 0.996 to 0.9996 after arm-phase calibration (0.96 at the two
extreme ports, where taper coupling, absent from this model, matters).
On the optimised design the BPM and the exact-path model give the same
field-averaged efficiency to 0.6\%.

\hypertarget{3-what-the-beam-propagation-found}{%
\section{What the beam propagation
found}\label{fpr_design-3-what-the-beam-propagation-found}}

(i) \textbf{Aberration of the confocal coupler.} With ports on an arc of
radius R, each off-axis port\textquotesingle s arm phases deviate from a
straight line by up to 1.3 rad pk-pk at the edge port (0.43 rad half-way
out), confirmed by exact path lengths. Arm trims cannot remove it
because it depends on the port. Placing the ports on an arc of radius
0.555 R (close to the Rowland circle) halves it over the field of use;
with 192 arms the worst residual is 0.08 rad rms. The Rowland-type arc
gives each port a different piston phase, which the static port trims of
route E remove (without them, adjacent ports in a window do not add in
phase).

(ii) \textbf{Coupling between dense ports.} In the converging 1.7 um
tapers (0.3 um gaps) neighbouring ports exchange light. With fixed
Gaussian windows this broadens the spot tails, and moving one of 20
tweezers then changes the others by up to 7 nm and 1.1\%. All of this is
linear and static: least-squares synthesis on the calibrated
port-to-atom response removes it (\textless= 3.5 nm, \textless= 0.12\%;
switch-off residual 1.3e-6) at 1 to 2\% cost in efficiency. Synthesis
also equalises depth over the field without a lookup. It is one matrix
product per update (precomputed 128 x n matrix).

\hypertarget{4-geometry-optimisation}{%
\section{Geometry optimisation}\label{fpr_design-4-geometry-optimisation}}

Figure of merit: the field-averaged laser power per tweezer for equal
depth, i.e. the harmonic mean of the peak intensity per unit port power
over \textbar x\textbar{} \textless= 60 um. Constraints: neighbour
intensity at 4 um \textless= 1e-4; gaps

\begin{quote}
= 0.3 um; 128 ports (fast actuators) fixed.
\end{quote}

The roll-off across the field has one physical origin on both sides of
the chip: scan loss. The arms capture tilted light from off-axis ports,
and the emitters radiate tilted beams, each with an element pattern that
falls across its angular range. Two levers act on it:

{\small
\begin{longtable}[]{@{}llll@{}}
\toprule\noalign{}
arms (128 ports) & facet 0.6 um: gain / edge ratio & facet 0.9 um &
facet 1.2 um \\
\midrule\noalign{}
\endhead
\bottomrule\noalign{}
\endlastfoot
128 & 1.01 / 0.53 & 1.10 / 0.51 & 1.23 / 0.46 \\
160 & 1.12 / 0.67 & 1.24 / 0.65 & 1.41 / 0.61 \\
192 & 1.18 / 0.76 & 1.31 / 0.74 & 1.51 / 0.71 \\
256 & 1.24 / 0.86 & 1.39 / 0.85 & 1.61 / 0.83 \\
\end{longtable}
}

(Gain relative to the confocal 128 x 128 design with 0.6 um facets;
exact-path model, Rowland-type arc.)

\textbf{Oversampling} (more arms and emitters than ports) confines the
ports to the central part of the arms\textquotesingle{} angular range
and raises the edge capture from 0.44 to 0.54 (192 arms). Arms and
emitters carry only static trims in route E, so this costs area and slow
heaters, not fast actuators. 192 arms take most of the benefit; 256 add
7\% for 64 more trims and a 1.8 mm slab.

\textbf{Wider emitter facets} narrow the element pattern and raise the
relay transmission at the centre from 0.64 (0.6 um) to 0.84 (1.2 um at
1.5 um pitch). At 0.5 um gaps the transmission saturates near 0.74 to
0.76 for any pitch from 1.5 to 2.2 um, so the 0.3 um gap of the 1.2 um
facet is what buys the last 10\%. Its coupling length is 227 um, which
limits the wide facet section to about 15 um.

\textbf{Slab widths.} The widest arms and ports (1.7 um at 2.0 um pitch)
win: field-averaged efficiency 0.261 against 0.18 to 0.24 for 1.2 to 1.5
um, because fill factor outweighs the slightly better angular acceptance
of narrow arms. The 0.3 um gaps are safe at the slab: the coupling
length is 622 um and the tapers converge over about 20 um.

\textbf{Window width.} Narrowing the Gaussian window from 0.9 to 0.85
ports gains 5\% but raises the neighbour intensity to 4.8e-5;
least-squares synthesis with a 1.68 um target is the adopted compromise.

\hypertarget{5-limits-and-open-points}{%
\section{Limits and open points}\label{fpr_design-5-limits-and-open-points}}

The model is 2D with a calibrated effective index; the residual 3D
effects are polarisation coupling at sidewalls and slab-edge scattering.
Fabrication variation of the 0.3 um gaps and of the taper geometry has
not been simulated; static errors are absorbed by the trims and the
measured response matrix, but the 1\% actuator accuracy for 10 nm
position noise (\texttt{m10\_route\_E\_LS\_*}, noise) remains. The relay
and objective are scalar and aberration-free. Steering is along one row.

Next steps: (i) tolerance analysis of gap and width errors in the dense
tapers; (ii) 3D FDTD of a few port-taper-slab cells for polarisation and
scattering; (iii) GDS layout of D1 with its 255 fast actuators (127 tree
splitters, 128 port phase shifters, which also carry the port trims) and
192 static arm trims; (iv) the calibration procedure that measures the
port-to-atom response on the device.

\hypertarget{verification-notes}{%
\section{Verification notes}\label{fpr_design-verification-notes}}

All numbers are SIM outputs of this repository except the loss-chain
assumptions marked above (CHOICE).

{\small
\begin{longtable}[]{@{}>{\raggedright\arraybackslash}p{0.320\dimexpr(\linewidth-8\tabcolsep)\relax}>{\raggedright\arraybackslash}p{0.094\dimexpr(\linewidth-8\tabcolsep)\relax}>{\raggedright\arraybackslash}p{0.086\dimexpr(\linewidth-8\tabcolsep)\relax}>{\raggedright\arraybackslash}p{0.500\dimexpr(\linewidth-8\tabcolsep)\relax}@{}}
\toprule\noalign{}
value & location & type & pointer \\
\midrule\noalign{}
\endhead
\bottomrule\noalign{}
\endlastfoot
calibrated indices, coupling lengths & Sec. 2 & SIM &
results/m10\_eim\_calibration.json \\
BPM convergence, BPM vs FDTD & Sec. 2 & SIM &
results/m10\_bpm\_small.json, m10\_bpm\_vs\_fdtd.json \\
confocal 128 x 128 capture and aberration & Sec. 3 & SIM &
results/m10\_fpr\_full\_A.json \\
optimised D1 capture and aberration & Sec. 1 & SIM &
results/m10\_fpr\_opt\_A.json \\
window and least-squares evaluations & Sec. 1, 3 & SIM &
results/m10\_route\_E\_\{full,opt\}\emph{A.json,
m10\_route\_E\_LS}\{full,opt\}\_A.json \\
optimisation tables & Sec. 4 & SIM & results/m11\_stage1.json,
m11\_stage2.json \\
power per tweezer & Sec. 1 & DERIVED & m9\_E\_power.json scaled by the
field-averaged peak per unit power \\
\end{longtable}
}
